\scanpage{26}
\deleted{Comme $T_p(u)$ et $T_p(u^*)$ sont}

\deleted{En effet, d'après la théorie p-adique \inserted{de Serre} ou th.
caractéristique de Weil, on aura
\[
(*)\qquad\deg u=\det T_p(u)_{\mathrm{\acute et}}\,\det T_p(u^*)_{\mathrm{\acute et}},
\]
et il reste à prouver que \unsure{$-=$|$=$}
\[
\begin{aligned}
1^\circ)\quad&T_p(u)_{\mathrm{\acute et}}=T_p(u^*)_{\mathrm{\acute et}}\\
2^\circ)\quad&\varepsilon(u)=1
\end{aligned}
\]
\inserted{ce qui implique $\operatorname{d\acute et}(T_p(u))=\varepsilon(u)\sqrt{\deg u}$, $\varepsilon(u)=\pm1$}}

\deleted{D'ailleurs on peut \inserted{se borner à} supposer
$u\in$ $\operatorname{Hom}(A,A)\otimes_{\mathbb Z}\mathbb Q\underset{\varepsilon}{=}\unsure{E|F}$,
et alors $\varepsilon$ est un hom $E^*\to(+1,-1)$, \deleted{\unsure{donc|de plus}}
qui est une fonction \unsure{rationnelle|rationelle} [quotient de
\deleted{\unsure{on voit aussitôt que tout tel hom est|assujettie que tout tel hom est}}
\deleted{identiqumt} deux fonctions}
% archive sheet number: 25
